
\section*{MATH311 Комплекс анализ хичээлийн ерөнхий мэдээлэл}

\begin{tabular}{ll}
Хичээлийн түвшин: & Бакалавр (Математикийн анги) \\
Кредит: & 3 кредит \\
Лекц: & 2 цаг / долоо хоног \\
Семинар: & 2 цаг / долоо хоног \\
Өмнөх мэдлэг: & Анализ I, II, Олон хувьсагчийн функцийн онол
\end{tabular}

%-------------------------------------------------

\section*{MATH311 Комплекс анализ хичээлийн тодорхойлолт}

Энэ хичээлээр комплекс тоо болон комплекс хувьсагчийн функцийн үндсэн онолыг судална. Комплекс функцийн хязгаар, тасралтгүй чанар, уламжлал, интеграл, цуваа, онцгой цэгүүд, үлдэгдэцийн теорем болон түүний хэрэглээг онол–нотолгоо, бодлого бодолттой хослуулан үзнэ. Хичээл нь цаашдын математик анализ, дифференциал тэгшитгэл, математик физикийн хичээлүүдийн суурь болно.

%-------------------------------------------------

\section*{Сургалтын үр дүн}

Хичээлийг амжилттай дүүргэсний дараа оюутан:
\begin{itemize}
\item Комплекс тооны алгебр ба геометр шинжийг тайлбарлаж чадна.
\item Комплекс функцийн хязгаар, тасралтгүй чанарыг шалгаж чадна.
\item Коши--Риманы тэгшитгэлийг ашиглан уламжлал орших эсэхйиг тогтооно.
\item Голоморф ба гармоник функцийн уялдааг ойлгоно.
\item Кошигийн интеграл теорем, томьёог хэрэглэнэ.
\item Тэйлор, Лораны цуваагаар функц задлан шинжилнэ.
\item Онцгой цэгүүдийг ангилж, үлдэц ашиглан интеграл бодно.
\end{itemize}

%-------------------------------------------------
\newpage
\section*{Лекцийн сэдвийн хөтөлбөр (16 долоо хоног)}

\begin{enumerate}[label=\textbf{Хичээл \arabic*.}]

\item Комплекс тоо ба комплекс хавтгай

\item $\mathbb{C}$-ийн топологийн суурь ойлголтууд  
(дотоод, гадаад, хязгаарын цэг, тусгаарлагдсан цэг, цэгийн орчин, задгай ба битүү олонлог,  холбоост ба тусгаарлагдсан олонлог, муж)

\item Комплекс функцийн жишээнүүд, олон утгатай функцүүд  
(олон гишүүнт, рационал, экспоненциал, тригонометр функцүүд)

\item Комплекс функцийн хязгаар ба тасралтгүй чанар

\item Комплекс уламжлал ба Коши--Риманы тэгшитгэл  
(зайлшгүй ба хүрэлцээтэй нөхцөл)

\item Голоморф ба гармоник функцүүд, дээд эрэмбийн уламжлал

\item Конформ дүрслэл ба уламжлалын геометр утга

\item Комплекс интеграл

\item Кошигийн интеграл теорем  

\item Кошигийн интеграл томьёо

\item Тэйлорын цуваа, аналитик ба гармоник функц, тэдгээрийн уялдаа

\item Лораны цуваа ба онцгой цэгүүд

\item Үлдэц (Residue)

\item Үлдэцийн теоремын хэрэглээ

\item Конформ дүрслэл ба Мёбиусын хувиргалт

\item Геометрийн зарчмууд (аргументын зарчим, хамгийн их модулийн зарчим, Шварцын лемма, Алгебрийн үндсэн теоремын баталгаа)
\end{enumerate}

%-------------------------------------------------
\newpage
\section*{Үнэлгээний журам (108 оноо)}

\subsection*{Ирц — 21 оноо}
\begin{itemize}
\item Лекцийн ирц: 15 оноо (1 тасалбал 1 оноо хасна)
\item Семинарын ирц: 6 оноо (1 тасалбал 1 оноо хасна)
\item Лекц эсвэл семинарыг \textbf{6 удаа тасалбал} \textbf{F*} үнэлгээ авна
\end{itemize}

\subsection*{Семинарын оноо  — 25 оноо}
Хичээлийн явцад семинарын хичээл дээр 5 удаа самбарт аман асуултад хариулна.\\
Нэг удаагийн асуулт: 5 оноо\\ Асуултын агуулга:
\begin{itemize}
\item Комплекс тооны тооцоолол, зэрэг, логарифм
\item Уламжлал, Коши--Риманы тэгшитгэл
\item Голоморф, гармоник функц, конформ дүрслэл
\item Кошигийн теорем ба томьёо
\item Цуваа, онцгой цэг, үлдэц
\end{itemize}

\subsection*{Гэрийн даалгавар — 7 оноо}

\subsection*{Явцын шалгалт — 25 оноо}
\begin{itemize}
\item 9-р долоо хоногт авна
\item Эхний 7 лекцийн агуулга (интегралаас өмнөх)
\end{itemize}

\subsection*{Улиралын шалгалт — 30 оноо}
\begin{itemize}
\item 17-р долоо хоногт авна
\item Сүүлийн 9 лекцийн агуулга (интегралаас хойших)
\end{itemize}

\subsection*{Нийт оноо}
108 онооноос авсан оноогоор дүн гаргана.  
100-аас дээш оноо авсан тохиолдолд 100 гэж дүгнэнэ.

%-------------------------------------------------

\section*{Сурах бичиг, материал}

\begin{itemize}
\item С.Чимэдцэрэн, С.Батболд — \emph{Нэг хувьсагчийн функцийн комплекс анализ}
\item С.Чимэдцэрэн — \emph{Комплекс хувьсагчийн функцийн онол}
\item L. Ahlfors — \emph{Complex Analysis}
\item J. B. Conway — \emph{Functions of One Complex Variable}
\item E. Stein, R. Shakarchi — \emph{Complex Analysis}
\end{itemize}

%-------------------------------------------------
